\section{Variational Diffusion Models：统一框架}

\subsection{一般形式}

Variational Diffusion Models 定义了一个通用的扩散模型家族。跳跃分布的一般形式为：

$$
q(\mathbf{x}_t \mid \mathbf{x}_0) = \mathcal{N}\!\left(\mathbf{x}_t;\, \alpha_t \mathbf{x}_0,\, \sigma_t^2 \mathbf{I}\right)
$$

其中 $\alpha_t$ 和 $\sigma_t$ 是关于时间 $t$ 的可微函数，满足：

\begin{enumerate}
    \item $\alpha_0 = 1$，$\sigma_0 = 0$（初始时刻为原始数据）；
    \item $\alpha_T \approx 0$，$\sigma_T \approx 1$（终止时刻近似标准高斯）；
    \item \textbf{信噪比 SNR$(t) = \alpha_t^2/\sigma_t^2$ 关于 $t$ 单调递减}。
\end{enumerate}

\begin{importantbox}{SNR 单调递减的含义}
SNR 单调递减意味着：随着前向过程的推进，信号成分（$\alpha_t \mathbf{x}_0$）相对于噪声成分（$\sigma_t \boldsymbol{\epsilon}$）持续减弱。这是扩散模型"加噪"的本质——信息逐渐被噪声淹没，最终数据分布退化为纯噪声分布。
\end{importantbox}

\subsection{SMLD 和 DDPM 作为特例}

\begin{table}[H]
\centering
\begin{tabular}{lccc}
\toprule
\textbf{模型} & $\alpha_t$ & $\sigma_t$ & \textbf{SDE 类型} \\
\midrule
SMLD / NCSN & $1$ & $\sigma(t)$ & VE-SDE \\
DDPM & $e^{-\frac{1}{2}\int_0^t \beta(s)\,\mathrm{d}s}$ & $\sqrt{1 - e^{-\int_0^t \beta(s)\,\mathrm{d}s}}$ & VP-SDE \\
\bottomrule
\end{tabular}
\caption{SMLD 和 DDPM 在 Variational Diffusion Models 框架下的参数}
\end{table}

\begin{knowledgebox}{设计空间的开放性}
VE-SDE 和 VP-SDE 只是两个具体例子。任何满足 SNR 单调递减条件的 $(\alpha_t, \sigma_t)$ 对都定义了一个合法的扩散模型。这为研究者提供了巨大的设计空间——可以通过选择不同的噪声调度来优化生成质量、采样速度或其他指标。例如，EDM（Elucidating the Design Space of Diffusion-Based Generative Models, Karras et al. 2022）就系统性地探索了这一设计空间。
\end{knowledgebox}

\subsection{从 SDE 到跳跃分布}

给定前向 SDE 的 $f(\mathbf{x}, t)$ 和 $g(t)$，可以推导出对应的 $\alpha_t$ 和 $\sigma_t$：

$$
\alpha_t = e^{\int_0^t f(s)\,\mathrm{d}s}, \qquad \sigma_t^2 = \alpha_t^2 \int_0^t \frac{g^2(s)}{\alpha_s^2}\,\mathrm{d}s
$$

反过来，给定 $\alpha_t$ 和 $\sigma_t$，也可以反推出 $f(t)$ 和 $g(t)$：

$$
f(t) = \frac{\mathrm{d} \ln \alpha_t}{\mathrm{d}t}, \qquad g^2(t) = \frac{\mathrm{d}\sigma_t^2}{\mathrm{d}t} - 2\frac{\mathrm{d}\ln \alpha_t}{\mathrm{d}t}\sigma_t^2
$$

这种双向转化意味着\textbf{三种等价描述}——转移分布、跳跃分布、SDE——可以自由互换。

\begin{warningbox}{符号混淆警告}
不同论文使用不同的符号约定。例如，Song et al. 的 $\sigma$ 指 VE-SDE 中的噪声调度（即 SMLD 的 $\sigma_i$），而 Variational Diffusion Models 中的 $\sigma_t$ 是跳跃分布的标准差。在同一个推导中可能同时出现两种 $\sigma$，需要根据上下文区分。Minhyuk Sung 教授在授课时也特别提醒了这一点。
\end{warningbox}

\subsection{本章小结}

Variational Diffusion Models 通过 $(\alpha_t, \sigma_t)$ 参数化提供了一个通用框架，SMLD 和 DDPM 是其特例。三种描述（转移分布、跳跃分布、SDE）之间存在精确的双向对应关系。


